\begingroup
\singlespacing
\begin{tikzpicture}[>=Stealth,
 token/.style={draw,minimum width=1.25cm,minimum height=0.8cm,font=\small},
 note/.style={align=center,font=\small}]
\node[note] at (2.5,1.15) {Entrada y contexto};
\node[note] at (7.5,1.15) {Respuesta de referencia};
\node[note] at (11.25,1.15) {Padding};
\foreach \k in {0,...,3} {
 \node[token,fill=gray!15] at (0.625+1.25*\k,0) {$x_{\number\numexpr\k+1\relax}$};
 \node[note] at (0.625+1.25*\k,-0.85) {$0$};
}
\foreach \k in {0,...,3} {
 \node[token,fill=blue!20] at (5.625+1.25*\k,0) {$y_{\number\numexpr\k+1\relax}$};
 \node[note] at (5.625+1.25*\k,-0.85) {$1$};
}
\foreach \k in {0,1} {
 \node[token,fill=white,dashed] at (10.625+1.25*\k,0) {pad};
 \node[note] at (10.625+1.25*\k,-0.85) {$0$};
}
\node[note,anchor=west] at (0,-1.55) {Máscara de pérdida $m_t$: $1$ incluye el token y $0$ lo excluye};
\draw[->,thick] (0,-2.1)--(10,-2.1);
\node[note] at (5,-2.6) {Cada token de respuesta se predice con los tokens anteriores};
\end{tikzpicture}
\endgroup
